\centerline{\bf \huge 10 класс}
\centerline{\normalsize{\it Работа рассчитана на \bf{180} минут}}

\problem{10.1}
Вася написал на доске несколько различных двузначных чисел. Оказалось, что сумма никаких двух из написанных на доске чисел не равна 90. Какое максимальное количество чисел мог написать Вася?

\problem{10.2}
Дан квадратный трёхчлен $f(x)$. Известно, что существуют пары различных чисел $(m, k)$ таких, что $\frac{f(m)}{f(k)}=\frac{m}{k}$. Назовем такие пары хорошими. Докажите, что значение выражения $m \cdot k$ во всех хороших парах одинаково.

\problem{10.3}
На доске были написаны не обязательно разные неотрицательные целые числа. Коля вычел из каждого исходного числа 1 , затем сложил модули всех получившихся чисел, и получил сумму $S_1$. Вася вычел из каждого исходного числа на доске 2 , затем сложил модули всех получившихся чисел, и получил сумму $S_2$ Наконец Андрей вычел из каждого исходного числа на доске 3 , затем сложил модули всех получившихся чисел, и получил сумму $S_3$. Сколько двоек было написано на доске?

\problem{10.4}
Есть 30 нечётных натуральных чисел. Оказалось, что их можно разбить на 10 троек так, что в каждой тройке одно из чисел равно произведению двух других. Может ли сумма этих 30 чисел равняться 2024?

\problem{10.5}
В трапеции $A B C D$ длина боковой стороны $A B$ равна сумме длин оснований. Окружность $\Omega$ с центром в точке $A$, проходящая через $D$, пересекает отрезки $B D$ и $B A$ в точках $E$ и $F$ соответственно. Докажите, что точки $B, C, E, F$ лежат на одной окружности.